\documentclass[11pt]{article}
\usepackage[margin=25mm]{geometry}
\usepackage{fontspec}
\setmainfont{Latin Modern Roman}
\usepackage{amsmath,amssymb,mathtools}
\usepackage{unicode-math}
\setmathfont{Latin Modern Math}
\usepackage{xurl,hyperref}
\hypersetup{colorlinks=true,urlcolor=blue}
\setcounter{secnumdepth}{0}
\setlength{\parindent}{0pt}
\setlength{\parskip}{5pt}
\setlength{\emergencystretch}{3em}
\newcommand{\sref}[2]{\hyperlink{source:#1:#2}{[S#2]}}
\newcommand{\cataloguescope}[1]{\paragraph{Recorded target.}#1}
\newcommand{\R}{\mathbb R}
\newcommand{\C}{\mathbb C}
\newcommand{\N}{\mathbb N}
\newcommand{\Z}{\mathbb Z}
\newcommand{\Q}{\mathbb Q}
\newcommand{\F}{\mathbb F}
\newcommand{\eps}{\varepsilon}
\begin{document}
% UnsolvedMath ID 20002549; source catalogue code AIM-PHYSICS-0057
\setcounter{section}{20002548}
\section{Zero-total-mass reduction and obstruction to the Euler integral}\label{20002549}
{\small UnsolvedMath ID 20002549 · Mathematical Physics}
\subsection{Definitions and mathematical statement}

\paragraph{Shared notation.}$\mathbb N=\{1,2,\ldots\}$; the standard number fields have their usual meanings. All additional conventions are stated locally.


\paragraph{Definitions and notation.}In the circular restricted three-body problem, two primaries follow prescribed circular motion and a test particle has negligible effect on them. In rotating coordinates the Jacobi constant combines kinetic energy and the effective gravitational/centrifugal potential. The Hill region is the set where the zero-velocity inequality permits motion. Its necks occur near collinear equilibria. A first integral is constant along trajectories; functional independence and Poisson commutation are needed for the usual Hamiltonian notion of complete integrability. The signed-mass condition $m_1+m_2=0$ is exceptional and cannot hold for two positive primaries. It must be interpreted in the generalized model intended by the source, not ordinary positive-mass circular dynamics.

\paragraph{Catalogue wording retained for traceability.}
In the circular restricted three-body problem, a Jacobi integral is available and the source says complete integrability requires two more integrals. It records an extra integral in the exceptional case $m_1+m_2=0$. Is there one more integral?
\paragraph{Scope and source qualification.}The signed-mass model, degree-of-freedom convention and admissible class of integrals need specification. Treating m1+m2=0 as compatible with the usual positive circular primaries would be inconsistent; the source question is preserved as an exceptional generalized setting. The catalogue title/status was paired with a machine-generated research attempt; that attempt is not peer-reviewed evidence of settlement. The statement covers the original source questions rather than only the title’s subcase.
\subsection{Short English statement}
Find whether the exceptional zero-total-mass restricted three-body model has the further first integral asked for in the source.
\subsection{Research attempt}
\paragraph{Model, date and investigated scope.}
OpenAI GPT-6 Astra (configured model \texttt{gpt-6-astra}), ultra reasoning; generated 11 October 2026. Derived zero-total-mass relative kinematics and a coaccelerating restricted Hamiltonian with its energy integral. The intended exceptional model and further integral remain unresolved.

\paragraph{Sources and literature check.}
Read the supplied complete definition, catalogue wording, scope qualifications and source list for this canonical ID. The original-source citations are inherited as supplied references; no fresh primary-source or novelty search was performed for this attempt. No current open-status claim is inferred from this calculation.

\paragraph{Testing the zero-total-mass kinematics.}In the signed two-primary Newtonian equations, let \(r=q_1-q_2\). Then
\[\ddot q_1=-m_2r/|r|^3,\quad \ddot q_2=m_1r/|r|^3,\quad
\ddot r=-(m_1+m_2)r/|r|^3.
\]
If \(m_1=m\ne0\) and \(m_2=-m\), the relative acceleration vanishes. A prescribed nonzero constant separation \(r_0\) is compatible with both primaries sharing the constant acceleration \(a=mr_0/|r_0|^3\). It is not compatible with an ordinary nonzero-frequency circular relative motion, whose relative acceleration would be \(-\omega^2r_0\).
\paragraph{The coaccelerating restricted equation.}Write the primary trajectories as \(Q(t)+a_i\), with \(Q''=a\), and the test-particle position as \(Q(t)+x(t)\). Its equation is
\[\ddot x=-\nabla U(x),\quad
U(x)=-\frac{m}{|x-a_1|}+\frac{m}{|x-a_2|}+a\cdot x.
\]
The signs follow by differentiating the two Coulomb terms and subtracting the common frame acceleration. Hence \(\tfrac12|\dot x|^2+U(x)\) is conserved on collision-free trajectories. The inertial-frame correction is a linear potential; simply replacing the signed pair by fixed centers and omitting it changes the model.
\paragraph{Gap.}The source's exceptional generalized setting must specify whether this coaccelerating model or some imposed rotating-primary system is meant, and whether there are two or three degrees of freedom. The calculation neither identifies the recorded extra integral nor constructs another independent Poisson-commuting integral. The first missing step is matching the precise Hamiltonian and known integral, then testing an admissible class of additional integrals. No claim of nonintegrability follows from the kinematic inconsistency of ordinary positive-mass circular language.

\paragraph{Verification and adversarial audit.}
Checked signed accelerations, vanishing relative acceleration, nonzero circular inconsistency, linear frame potential sign and lack of a nonintegrability claim. This was an LLM adversarial audit, not independent human verification; no theorem prover was used. Prompts A then B applied to this ID: Checked signed accelerations, vanishing relative acceleration, nonzero circular inconsistency, linear frame potential sign and lack of a nonintegrability claim. Scope remains limited: Identify the precise generalized Hamiltonian, degrees of freedom and recorded extra integral, then analyze independent commuting integrals.

Coordinator checked the signed-mass relative acceleration and the linear inertial potential in the translating accelerated frame. This additional check was model-assisted, not independent human verification.

\paragraph{Tools, process and contributions.}
Python 3 and local shell used to read supplied files and serialize this manuscript; mathematics derived symbolically in the text, with no numerical experiment or formal proof assistant. Contribution prompt A was followed by prompt B. GPT-6 Astra with ultra reasoning drafted the research and performed its stated self-audit; any additional coordinator contribution is identified in the verification and audit. A GPT-6-based Codex coordinator assembled the manuscript and checked repository integration; its exact variant and reasoning configuration were not exposed. No human mathematical contribution, independent human proof check or formal proof certification is supplied. Complete derivations and the consulted source references are included here; private reasoning transcripts are not published.

\paragraph{Final status and first remaining gap.}
\textbf{UNRESOLVED}. Identify the precise generalized Hamiltonian, degrees of freedom and recorded extra integral, then analyze independent commuting integrals.

\section{Completion Estimate}
\noindent\textbf{COMPLETION: 0\%}
The calculation is an elementary benchmark or conditional reduction; it establishes no claimed new portion of the complete source target.

\subsection{Sources}
\begin{itemize}
\item[\hypertarget{source:20002549:1}{S1}] ULAM.AI and the UnsolvedMath Contributors, UnsolvedMath, record 20002549 (AIM-PHYSICS-0057). Catalogue curation under CC BY 4.0; definitions and English exposition added locally.
\url{https://huggingface.co/datasets/ulamai/UnsolvedMath}.
\item[\hypertarget{source:20002549:2}{S2}] American Institute of Mathematics, Variational Methods in Celestial Mechanics; notes by Kuo-Chang Chen, version 22 June 2003, source item 20.
\url{https://aimath.org/WWN/varcelest/varcelest.pdf}.
\end{itemize}
\label{end:20002549}
\end{document}
